The strongest patient-level discrimination in this sweep emerged from the sparsest tested configuration. The best hybrid run used \texttt{inner-k=1}, selected only 7 unique features appearing in at least one fold, and exceeded the best linear baseline by a large margin on ROC-AUC and MCC. Figure~\ref{fig:sweep-innerk-roc} also shows that this was not a generic effect of smaller feature budgets across both models. In the low-budget region, the hybrid model retained stronger ROC-AUC than the linear SVM and reached its peak performance at the sparsest tested setting, whereas the linear model reached its best ROC-AUC only after a broader per-fold budget. In this cohort, the value of the more expressive model and the value of strong sparsity appear to have operated together.

Statistical inference was assessed with permutation tests and three tiers of multiple-comparisons correction to account for the structured sweep design. The raw permutation p-value of $p=0.008$ survives the primary structured correction (2 model families, $p_\text{corr}=0.016$) at $\alpha=0.05$. As robustness checks, we computed a max-T permutation correction across all 150 configurations ($p_\text{corr}=0.302$) and a full-sweep Bonferroni ($p_\text{corr}=1.00$), the latter being mechanically uninformative given the 150-test multiplier. We adopt the 2-family correction as the primary inference because the inner-k and window-size factors represent exploratory levels of a single hypothesis rather than independent comparisons, and report all tiers for complete transparency. Several additional caveats qualify the results. The study is limited to $n=21$ participants with a 7:14 class imbalance, producing wide bootstrap confidence intervals (e.g., ROC-AUC 95\% CI: 0.600--0.990) that span both chance-level and near-perfect performance. The hybrid model's hyperparameters were inherited from a standard configuration rather than tuned specifically on this cohort; while this reduces researcher degrees of freedom, the configuration may not be optimal for the present data.

One plausible explanation for the performance gap between models is that aggressive per-fold pruning removed weak or redundant handcrafted features and exposed a compact signal set whose interactions were more useful to the hybrid architecture than to a linear decision boundary. Supporting this interpretation, eight classical classifiers and two neural baselines (2-layer MLP, GAP architecture variant) were evaluated at the identical \texttt{inner-k=1} configuration; all performed at or below chance (ROC-AUC 0.204--0.296, Table~\ref{tab:intermediate-baselines}). An additional ablation removed SMOTE from the hybrid, reducing ROC-AUC from 0.816 to 0.582. Together these results show that the full CNN-BiLSTM-attention pipeline with SMOTE-based training is jointly required — neither non-linearity alone, nor SMOTE alone, nor a simplified deep architecture suffices. With only a few retained descriptors, the convolutional, recurrent, and attention stages appear better positioned to combine complementary temporal and frontal-temporal information without dilution from a larger pool of unstable candidates. We did not compute formal paired statistical tests (e.g., DeLong or bootstrapped AUC comparison) between the two model families because the best configurations operate at different inner-k values (1 vs.\ 30). At matched inner-k values near the hybrid's optimum, the SVM operates near chance (ROC-AUC $\approx$ 0.4--0.5; Figure~\ref{fig:sweep-innerk-roc}), so the performance gap is large enough that formal testing at the matched optimum would add limited insight. The intermediate-baseline experiment (Table~\ref{tab:intermediate-baselines}) provides complementary evidence: ten alternative models all fail at inner-k=1, supporting the conclusion that the hybrid's advantage is not a statistical artifact of the sweep. This interpretation remains provisional, but it is consistent with the sweep pattern and with the marked separation between the best sparse hybrid run and the broader-budget linear baseline.

This result fits part of the EEG depression literature and supports a measured interpretation. Prior work has argued that EEG can provide clinically useful baseline markers in major depressive disorder and treatment-response settings \citep{olbrich2013_eeg,olbrich2015_personalized,leuchter2010_biomarkers}. Related discussions have also emphasized tighter biomarker discipline, better standardization, and more realistic expectations about generalization \citep{Fu2019,simmatis2023_technical,Moncy2025BD}. Machine learning studies have likewise reported encouraging treatment-response signals, including combinations of EEG and clinical covariates \citep{AlKaysi2017,jaworska2019_leveraging,oakley2022_eeg}. Other modelling studies reach similar conclusions while still exposing limits from data heterogeneity, cohort size, or insufficient interpretability \citep{shalbaf2018_nonlinear,ebrahimzadeh2021_predicting,zhdanov_svm_prediction}. Recent predictive overviews and reviews continue to highlight the same translational tension \citep{watts2022predicting,cohen2023electroencephalography}. Table~\ref{tab:sota-comparison} places the present results in the context of comparable EEG-based treatment-response prediction studies. The achieved ROC-AUC of 0.816 is competitive with prior work using larger montages and cohorts, despite the present study's substantially reduced channel count and sparse feature budget. However, direct numerical comparison is complicated by differences in clinical populations, outcome definitions, and validation methodology; the table is provided for contextual orientation rather than formal benchmarking.

\begin{paperwidetable}
  \centering
  \caption{Comparison with selected EEG-based treatment-response prediction studies in depression.}
  \label{tab:sota-comparison}
  \footnotesize
  \begin{tabular}{lcccccl}
\toprule
Study & Method & EEG channels & $n$ & AUC & Ext.\ val. & Task \\
\midrule
\citet{Moncy2024} & Deep learning (PSD features) & 4 (AF7,AF8,TP9,TP10) & 21 & 0.76$^\dagger$ & No & tDCS remission prediction \\
\citet{Xiao2025PMIP} & PLV connectivity + ML & 4 (AF7,AF8,TP9,TP10) & 21 & 0.72$^\dagger$ & No & tDCS remission prediction \\
\citet{AlKaysi2017} & SVM, RF, LR (clinical + EEG) & 19 & 80 & 0.82 & No & Antidepressant response \\
\citet{jaworska2019_leveraging} & SVM, LR (EEG + clinical) & 32 & 73 & 0.83 & No & Antidepressant response \\
\citet{oakley2022_eeg} & LORETA + SVM & 72 & 65 & 0.75 & No & Sertraline/placebo response \\
\citet{ebrahimzadeh2021_predicting} & Nonlinear EEG + classifier & 19 & 30 & 0.90 & No & Antidepressant response \\
\citet{shalbaf2018_nonlinear} & Nonlinear features + SVM & 19 & 45 & 0.88 & No & rTMS response \\
\citet{zhdanov_svm_prediction} & SVM with spectral features & 19 & 44 & 0.76 & No & Treatment response \\
\textbf{This study} & \textbf{Sparse CNN-LSTM + handcrafted features} & \textbf{4 (AF7,AF8,TP9,TP10)} & \textbf{21} & \textbf{0.82} & \textbf{No} & \textbf{tDCS remission prediction} \\
\bottomrule
\end{tabular}

\smallskip
\footnotesize $^\dagger$Estimated from reported metrics. AUC = area under the ROC curve (patient-level where available). Ext.\ val. = external validation cohort. Studies selected for relevance to EEG-based treatment response prediction in depression; the list is illustrative rather than exhaustive.

\end{paperwidetable}

The present sweep adds a feature-selection view in which larger feature budgets yield more overlap and many more unique selected features. Those larger budgets do not improve discrimination. The panels support a small set of frequently selected features with consistent direction of class separation. That reading keeps the modest Jaccard and Kuncheva scores in view and gives appropriate weight to the stronger effect-direction signal. However, the low Jaccard (0.190) and Kuncheva (0.187) scores indicate that the selected feature changes in most folds, which is a limitation for biomarker reproducibility claims. Three of the seven features appear in only a single fold, and the ``perfect'' sign consistency values for those entries should be interpreted cautiously. The Jaccard increase with larger inner-k is partly mechanical (selecting $k$ features from a fixed pool of 247 will produce more overlap by construction as $k$ increases), so this trend should not be over-interpreted as evidence of underlying signal stability. Methodologically, this framing matters because the depression EEG biomarker literature still contains substantial publication bias, limited out-of-sample validation, and known vulnerability to leakage from pre-cross-validation feature processing \citep{Widge2019,Shim2021}. The present pipeline addresses that problem by embedding selection and rebalancing inside participant-level leave-one-participant-out folds and then summarizing recurrence and directionality instead of reporting only a single global feature list \citep{Shen2025}.

The same clinical program has also supported other signal representations in recent PSD-based and PLV-based studies, which located informative response-related signal in low-density temporal and frontal-temporal patterns \citep{Moncy2024,Xiao2025PMIP}. In the present sweep, the most frequently selected features again concentrated in TP10-derived and frontal-temporal descriptors. That regional convergence supports the view that the low-montage recordings may contain response-relevant information across multiple feature families. However, the current study provides only within-cohort evidence; the generalizability of these features under different montage configurations, clinical populations, and treatment protocols remains untested.

The specific features selected in the present sweep are neurophysiologically plausible. Elevated temporal spectral entropy in non-remission participants suggests greater irregularity or complexity in the resting-state signal at the TP10 site, which overlies the temporoparietal junction — a region implicated in the default mode network and self-referential processing, both of which are dysregulated in depression. Higher entropy may reflect less organized or more variable neural dynamics that are less amenable to tDCS-induced modulation. The frontal-temporal beta and gamma difference features point toward inter-regional communication in higher frequency bands, which have been linked to top-down cognitive control and attentional processes mediated by fronto-temporal networks. tDCS applied over the dorsolateral prefrontal cortex (approximated by AF7/AF8) is thought to modulate excitability in these circuits, and baseline differences in frontal-temporal coupling may index a network's capacity to respond to such modulation. The AF7 zero-crossing count — a time-domain measure of signal volatility — showed the clearest remission-associated pattern and may capture a more dynamic or responsive frontal signal state that is conducive to stimulation-induced plasticity. These interpretations are necessarily speculative given the study's sample size and lack of independent replication, but they provide concrete, testable hypotheses for future mechanistic work.

Related tDCS studies in bipolar depression from the same group point in a similar low-montage direction. PSD-based modeling again highlighted AF7 and TP10 signal content, and PLV analyses emphasized frontal-temporal and temporoparietal synchronization patterns \citep{Moncy2025BD,Xiao2025JAD}. These bipolar studies address a different diagnosis and treatment protocol. They support reduced-montage feasibility and signal plausibility in adjacent affective cohorts.

This study has several limitations. First, the relatively small participant size ($n=21$) with a 7:14 class imbalance limits statistical power and generalizability; the wide bootstrap confidence intervals (e.g., ROC-AUC 95\% CI: 0.600--0.990) reflect this uncertainty. Although SMOTE and group equalization were applied within training folds, the fundamental class imbalance remains a constraint. Second, the primary corrected p-value ($p_\text{corr}=0.016$) survives the structured 2-family Bonferroni correction but does not survive the more conservative max-T correction across all 150 configurations ($p_\text{corr}=0.302$); replication in an independent cohort is needed. Third, there is no external validation cohort, so the reported performance and identified features cannot yet be assessed for generalizability to independent populations, alternative low-density montages, or different clinical settings. Fourth, important clinical variables such as medication status, illness duration, episode count, and psychiatric comorbidities were not available; these may influence both EEG features and treatment outcomes. Fifth, the 4-channel low-density montage trades spatial coverage for practical deployability. Recording only AF7, AF8, TP9, and TP10 captures frontal and temporal regions while omitting central, parietal, and occipital sites. This limits source-localization precision and precludes analysis of well-studied posterior alpha asymmetry markers. However, the reduced setup also enables a fully home-based, self-administered protocol that would be impractical with a 19- or 32-channel research montage, and the observed signal in these four channels supports the feasibility of this trade-off for remote tDCS monitoring. Whether a richer montage would yield higher predictive performance is an empirical question for future head-to-head studies. Sixth, the feature stability metrics (mean Jaccard 0.190, Kuncheva 0.187) indicate limited cross-fold consistency; however, median-based patient aggregation improved ROC-AUC to 0.857, suggesting that the feature signal is stable even when fold-to-fold feature identity varies. Seventh, when \texttt{inner-k=1}, SMOTE operates in a 1-dimensional feature space where the standard rationale for synthetic minority oversampling does not fully apply. The no-SMOTE ablation confirmed this concern: removing SMOTE reduced the hybrid's ROC-AUC from 0.816 to 0.582, so SMOTE is empirically essential at this setting despite the theoretical caveat. Eighth, the study does not benchmark against recent transformer-based or graph-neural-network EEG classifiers, which represent an active area of development. Ninth, the hybrid architecture contains approximately 1.2 million parameters and requires GPU acceleration for training; inference time and memory footprint, while modest for a research server, may constrain deployment on edge devices. Tenth, the SOTA comparison (Table~\ref{tab:sota-comparison}) is contextual rather than a formal benchmark due to differences in clinical populations, outcome definitions, and validation methodology between studies; direct method-level comparisons would require harmonized datasets and protocols.

The recurring features can be treated as plausible candidate features from this cohort that now warrant follow-up validation rather than as validated clinical biomarkers. Future work should test whether the same sparse feature set recurs under independent cohorts, alternative low-density montages, and explicit external validation \citep{Dagnino2023}. Adjacent response-prediction work in depression and neuromodulation provides useful targets for that next stage \citep{funk2008_prefrontal,bailey2023_concurrent,guo2024_restingstate}. Home-based and transport-oriented studies offer one clear direction for follow-up \citep{romeromarn2026_remotely,kratter2026_stanford,waller2026_transient}. More recent response-signature studies also provide relevant benchmarks for that next phase \citep{sheen2024_response,zeidabadi2025_prediction,reissmann2026_thetacordance}. Related protocol and review work in treatment-resistant depression can help frame that validation effort \citep{ulrich2025_alterations,pettorruso2024_overcoming}. These extensions should keep stringent patient-level evaluation and sparse interpretable feature budgets in the core study design, ideally with pre-registered analysis plans to limit researcher degrees of freedom.
